\documentclass[11pt]{article}

% --- Packages ---
\usepackage[utf8]{inputenc}
\usepackage[margin=1in]{geometry}
\usepackage{amsmath}
\usepackage{amssymb}
\usepackage{amsthm}
\usepackage{mathtools}
\usepackage{booktabs}
\usepackage{xcolor}
\usepackage{enumitem}
\usepackage{hyperref}
\hypersetup{colorlinks=true, linkcolor=blue!50!black, urlcolor=blue!50!black}

% --- Notation macros ---
\newcommand{\ind}[1]{\mathbf{1}\!\left[#1\right]}        % indicator
\newcommand{\Terms}{\mathcal{T}}                          % term set of a doc
\newcommand{\Vocab}{\mathcal{V}}                          % term vocabulary
\newcommand{\Base}{\mathcal{V}_{B}}                       % baseline vocabulary
\newcommand{\Ent}{\mathcal{E}}                            % NER entity vocabulary
\newcommand{\Wset}{\mathcal{W}}                           % discovery weeks
\newcommand{\Gw}{G_{w}}                                   % docs in week w
\renewcommand{\deg}{\delta}                               % reach / degree
\newcommand{\clu}{\gamma}                                 % clustering
\newcommand{\Rtab}{\mathsf{R}}
\newcommand{\Ptab}{\mathsf{P}}

\theoremstyle{definition}
\newtheorem{definition}{Definition}
\theoremstyle{remark}
\newtheorem{remark}{Remark}

\title{\textbf{Entity Emergence Detection}\\[2pt]
\large A graph-topological pipeline for early thematic detection from news headlines}
\author{Thematic Trading research}
\date{Notebook \texttt{0.20-genai-entity-emergence-ner}}

\begin{document}
\maketitle

\begin{abstract}
We formalise the pipeline that turns a stream of dated news headlines into a short list
of candidate investment \emph{themes}, and derive the mathematical reasons each stage
works. The pipeline is a \emph{high-recall-then-high-precision cascade}: a novelty filter
isolates newly emerging entity names, a weekly co-occurrence graph scores each entity by
two structural quantities (its \emph{reach} and its \emph{clustering}), a temporal rule
keeps only entities that persist, grow, and keep bridging unrelated stories, and a final
language-model pass removes the residual non-themes. The central idea is that a
\emph{theme} and an \emph{event} have opposite graph topologies: a theme is a low-clustering
hub that brokers across disconnected communities, an event is a high-clustering clique.
On the generative-AI benchmark the pipeline catches \texttt{chatgpt} in the week of
3--9 January 2023 and promotes it the week of 17--23 January 2023, roughly four months
ahead of the corresponding thematic ETF inception.
\end{abstract}

% =====================================================================
\section{Data model and notation}
% =====================================================================

The input is a corpus of headlines. Each headline (document) $d$ carries a timestamp
$\tau(d)$ and a set of multi-word \emph{terms} $\Terms(d)\subseteq\Vocab$, extracted upstream
(noun-phrase / key-term extraction). We write $\mathrm{tok}(t)$ for the word tokens of a
term $t$.

\paragraph{Time windows.}
Fix a baseline cutoff $t_0$ and a discovery interval $[t_1,t_2]$. In the notebook,
\[
t_0 = \text{2022-09-30},\qquad [t_1,t_2] = [\text{2022-10-01},\,\text{2023-02-28}].
\]
The baseline window is $B=\{d:\tau(d)\le t_0\}$. Time is bucketed into ISO weeks
($\texttt{W-MON}$); let $\Wset$ be the set of discovery weeks and $\Gw=\{d:\,\text{week}(d)=w\}$
the documents in week $w$.

\paragraph{Per-week statistics.}
For a term $t$ and week $w$ define the \emph{support} (mention count) and, for a pair $(a,b)$,
the \emph{co-occurrence count}
\begin{equation}
c_w(t)=\sum_{d\in\Gw}\ind{t\in\Terms(d)},\qquad
A_w(a,b)=\sum_{d\in\Gw}\ind{a\in\Terms(d)\ \wedge\ b\in\Terms(d)},\quad a\ne b.
\label{eq:counts}
\end{equation}
$A_w$ is the (symmetric, weighted) adjacency of the \emph{weekly co-occurrence graph}
$\mathcal{G}_w=(\Vocab_w,\,\{(a,b):A_w(a,b)>0\})$, whose nodes are the terms occurring that
week and whose edges join terms that share a headline.

\paragraph{Two tables flow through the pipeline.}
\begin{itemize}[leftmargin=1.4em]
  \item $\Rtab$, the \emph{detection table}: one row per $(\text{novel entity},\,\text{week})$,
        carrying $c_w$, reach $\deg_w$, clustering $\clu_w$, the local subgraph, and
        representative headlines.
  \item $\Ptab$, the \emph{per-entity lifecycle}: one row per caught entity, collapsing its
        weeks into \texttt{first\_caught}, \texttt{n\_weeks}, $\deg_{\max}$, median clustering,
        and a \texttt{promoted\_week}.
\end{itemize}

% =====================================================================
\section{Pipeline overview}
% =====================================================================

The pipeline is a sequence of monotonically shrinking sets. With the notebook's parameters
the funnel is
\[
\underbrace{3{,}767}_{\text{novel entities}}\ \longrightarrow\
\underbrace{814}_{\text{caught}}\ \longrightarrow\
\underbrace{20}_{\text{promoted}}\ \longrightarrow\
\underbrace{18}_{\text{themes}}\ \longrightarrow\
\underbrace{4}_{\text{LLM-kept}}.
\]
The stages are:
\begin{description}[leftmargin=2.2em,style=nextline]
  \item[Stage 0 -- entity selection.] Keep terms that name an entity (\S\ref{sec:stage0}).
  \item[Stage 1 -- novelty + support.] Keep terms absent from the baseline and seen
        $\ge m$ times that week (\S\ref{sec:stage1}).
  \item[Stage 2 -- graph features.] Score each surviving entity by reach $\deg_w$ and
        clustering $\clu_w$ (\S\ref{sec:stage2}). These three stages build $\Rtab$.
  \item[Stage 3 -- catch then promote.] Flag entity-weeks with enough reach
        (\emph{catch}), then keep entities that persist, grow, and stay bridges
        (\emph{promote}). This builds $\Ptab$ (\S\ref{sec:stage3}).
  \item[Stage 4 -- assemble.] Merge co-occurring promoted entities into themes
        (\S\ref{sec:stage4}).
  \item[Stage 5 -- LLM reject.] A conservative language-model filter removes clusters that
        are clearly not themes (\S\ref{sec:stage5}).
\end{description}

% =====================================================================
\section{Stage 0: entity selection}\label{sec:stage0}
% =====================================================================

\paragraph{Term cleaning.}
A term $t$ is admissible if it is long enough and contains no stop / calendar / generic-finance
token:
\begin{equation}
\mathrm{keep}(t)=\ind{|t|\ge 3}\;\wedge\;\ind{\mathrm{tok}(t)\cap \mathrm{STOP}=\varnothing}.
\end{equation}

\paragraph{Entity indicator.}
Let $\Ent$ be a \emph{corpus-wide entity vocabulary} of word tokens. A term is an entity iff
any of its tokens belongs to $\Ent$:
\begin{equation}
E(t)=\ind{\exists\,u\in\mathrm{tok}(t):\ u\in\Ent}.
\label{eq:isentity}
\end{equation}

\paragraph{Building $\Ent$.}
Run a named-entity recogniser (spaCy \texttt{en\_core\_web\_sm}) over all discovery
headlines. Collect the tokens appearing inside spans labelled in
\[
\mathcal{L}=\{\text{ORG},\text{PRODUCT},\text{PERSON},\text{WORK\_OF\_ART},\text{FAC},\text{NORP},\text{EVENT}\},
\]
and retain a token $u$ iff
\begin{equation}
u\in\Ent \iff
\bigl(\textstyle\sum_{\text{spans }s\in\mathcal{L}} \ind{u\in s}\ge 2\bigr)\ \wedge\
\bigl(u\notin \text{EnglishDict}\bigr)\ \wedge\ \bigl(u\ \text{matches}\ \texttt{[a-z][a-z0-9'\textbackslash-]+}\bigr).
\label{eq:vocab}
\end{equation}

\paragraph{Why a corpus-wide, non-dictionary vocabulary?}
Out-of-the-box NER is unreliable on Bloomberg-style Title-Case / ALL-CAPS headlines: it
mis-tags leading verbs as proper nouns (\emph{``Microsoft Mulls''} $\to$ \textsc{org}) and
misses lowercased product names (\texttt{ChatGPT} is dropped from $\approx 83\%$ of its
headlines). Two design choices repair this.
\begin{enumerate}[leftmargin=1.6em]
  \item \textbf{Aggregate, do not trust per-headline tags.} A token is admitted only if it is
        tagged as an entity in $\ge 2$ \emph{different} headlines. Casing noise is roughly
        independent across headlines, so a spurious tag (\emph{``Mulls''}) rarely recurs,
        whereas a true name (\emph{``openai''}) recurs. The count threshold is a denoiser:
        if $p$ is the per-headline false-tag probability for a junk token, the probability it
        clears the $\ge\!2$ bar across its occurrences falls like $O(p^2)$.
  \item \textbf{Keep only non-dictionary tokens.} Real entity names (\texttt{openai},
        \texttt{nvidia}, \texttt{chatgpt}) are almost never English dictionary words, whereas
        mis-tagged common words (\emph{mulls}, \emph{sparks}, \emph{creator}) are. Subtracting
        the dictionary, $\Ent \subseteq \Vocab\setminus\text{EnglishDict}$, removes the dominant
        NER error class while costing few true names.
\end{enumerate}
This converts the \emph{select-out} blocklist of the predecessor (an explicit list of
event verbs) into a \emph{select-in} whitelist of entity tokens learned from the corpus, which
generalises beyond a hand-curated list while agreeing with it on the cases that matter
(\texttt{chatgpt}, \texttt{openai}, \texttt{microsoft}, \texttt{nvidia} all pass both).

% =====================================================================
\section{Stage 1: novelty and support}\label{sec:stage1}
% =====================================================================

\paragraph{Baseline vocabulary.}
$\Base=\bigcup_{d\in B}\Terms(d)$ is the set of every term seen on or before $t_0$.

\paragraph{Novelty.}
A term is novel at week $w$ iff it never appeared in the baseline,
\begin{equation}
\mathrm{novel}(t)=\ind{t\notin\Base}.
\end{equation}
Set-difference against $\Base$ is a \emph{high-pass filter on term history}: it annihilates
every persistent, background, or recurrent name and keeps only first-time arrivals. Formally,
if $T(t)=\min\{\tau(d):t\in\Terms(d)\}$ is the term's first-appearance time, then
$\mathrm{novel}(t)=\ind{T(t)>t_0}$, so novelty is exactly a \emph{first-passage} (birth)
event relative to the baseline. This is what makes the detector forward-looking rather than
descriptive: emergence is, by construction, the appearance of vocabulary that did not exist
before.

\paragraph{Support.}
Among novel entities, keep only those with enough weekly mentions, $c_w(t)\ge m$ with $m=3$.
Under a null in which a junk term appears in headlines as independent rare events with weekly
rate $\lambda$, $c_w(t)\sim\mathrm{Poisson}(\lambda)$ and the per-week false-alarm rate is
$\Pr[c_w\ge 3]=1-e^{-\lambda}(1+\lambda+\tfrac{\lambda^2}{2})=O(\lambda^3)$ for small
$\lambda$. The cubic suppression of singletons and doubletons is why $m=3$ removes the long
tail of one-off names cheaply.

\paragraph{Candidate set.}
Combining Stages 0--1, the candidates at week $w$ are
\begin{equation}
C_w=\bigl\{t:\ E(t)=1,\ \mathrm{keep}(t)=1,\ t\notin\Base,\ c_w(t)\ge m\bigr\}.
\end{equation}

% =====================================================================
\section{Stage 2: graph features -- reach and clustering}\label{sec:stage2}
% =====================================================================

This is the heart of the method. For each candidate anchor $t\in C_w$ we read two numbers off
its neighbourhood in $\mathcal{G}_w$.

\paragraph{Entity neighbourhood and reach.}
Restrict $t$'s neighbours to other entities and order them by co-occurrence strength,
\begin{equation}
N_E(t)=\bigl\{p:\ A_w(t,p)>0,\ E(p)=1\bigr\},\qquad
\deg_w(t)=\bigl|N_E(t)\bigr|.
\end{equation}
$\deg_w(t)$ is the \emph{reach}: the number of distinct entity partners with which $t$ shares
a headline that week.

\paragraph{Local clustering of the partners.}
Take the $K$ strongest partners $P=\mathrm{top}_K\,N_E(t)$ (with $K=15$) and measure the edge
density \emph{among the partners themselves}:
\begin{equation}
\clu_w(t)=
\begin{dcases}
\frac{\displaystyle\sum_{\{i,j\}\subseteq P} \ind{A_w(P_i,P_j)>0}}{\dbinom{|P|}{2}}, & |P|\ge 2,\\[2ex]
0, & |P|<2.
\end{dcases}
\label{eq:clustering}
\end{equation}
Equation~\eqref{eq:clustering} is the \emph{local clustering coefficient} of node $t$
(restricted to its $K$ strongest entity partners): the fraction of partner pairs that are
themselves joined by an edge.

\paragraph{The theme--event dichotomy.}
The pair $(\deg_w,\clu_w)$ separates two qualitatively different ego-network topologies.
\begin{itemize}[leftmargin=1.6em]
  \item \textbf{Hub / star (theme).} $t$ links many partners that do \emph{not} link each
        other: $\deg_w$ large, $\clu_w\approx 0$. One concept (\texttt{chatgpt}) connects
        otherwise-unrelated firms and sectors. In Burt's language, $t$ \emph{spans structural
        holes}: information and exposure flow \emph{through} it because the communities it
        joins are not directly connected. This is the structural signature of a
        cross-sectoral theme.
  \item \textbf{Clique (event).} $t$ and its partners are all mutually connected: $\clu_w\approx 1$.
        A single story (\emph{FTX collapse}) makes the same fixed cast of actors co-occur in
        every headline, so the partner subgraph is near-complete. This is the signature of an
        idiosyncratic event, not a theme.
\end{itemize}

\paragraph{A generative justification.}
Suppose $t$'s partners are drawn from $C$ latent communities and two partners share a headline
with probability $p_{\mathrm{in}}$ if they belong to the same community and $p_{\mathrm{out}}$
otherwise, with $p_{\mathrm{out}}\ll p_{\mathrm{in}}$ (a stochastic-block / planted-partition
model). If $t$ is a broker drawing roughly evenly from all $C$ communities, the expected
clustering of its partner set is
\begin{equation}
\mathbb{E}\bigl[\clu_w(t)\bigr]\;\approx\;
\frac{1}{C}\,p_{\mathrm{in}}+\Bigl(1-\frac{1}{C}\Bigr)\,p_{\mathrm{out}}
\;\xrightarrow[\;C\ \text{large}\;]{}\;p_{\mathrm{out}}\approx 0 .
\end{equation}
A theme spans many communities ($C$ large) and therefore has clustering near the
\emph{between}-community edge probability, i.e. near zero. An event lives inside a single
community ($C=1$), giving clustering $\approx p_{\mathrm{in}}\to 1$. Thus $\clu_w$ is a direct,
unsupervised estimator of ``how many distinct narratives does this entity tie together'',
and low $\clu_w$ is precisely the cross-sectoral property a theme must have.

% =====================================================================
\section{Stage 3: catch and promote}\label{sec:stage3}
% =====================================================================

\subsection{Catch (high recall)}
An entity-week is \emph{caught} when it has both enough support and enough reach:
\begin{equation}
\mathrm{caught}(t,w)=\ind{c_w(t)\ge m}\ \wedge\ \ind{\deg_w(t)\ge \deg_{\min}},
\qquad m=3,\ \deg_{\min}=8.
\end{equation}
The reach bar $\deg_{\min}=8$ selects \emph{hubs}: it is a threshold on the horizontal axis of
the $(\deg,\clu)$ plane and removes small, parochial co-occurrence patterns regardless of their
clustering. Catch is deliberately permissive (it does not yet use $\clu$) so that recall stays
high; precision is recovered downstream.

\subsection{Promote (persist + grow + bridge)}
Collapse each caught anchor $a$ into its chronologically ordered caught weeks
$w_1<\dots<w_n$ with reaches $d_i=\deg_{w_i}(a)$ and clusterings $\clu_i=\clu_{w_i}(a)$.
The anchor is \emph{promoted} at the first index $i$ satisfying all three of
\begin{align}
&\textbf{persistence:} && i\ge \rho, &&\rho=2, \label{eq:persist}\\
&\textbf{growth:}      && d_i\ \ge\ \max_{j<i} d_j, &&\text{(record reach)} \label{eq:grow}\\
&\textbf{bridge:}      && \operatorname{median}(\clu_1,\dots,\clu_i)\ \le\ \kappa, &&\kappa=0.60. \label{eq:bridge}
\end{align}
The promoted week is that first $w_i$; if no week qualifies the anchor is never promoted
($\texttt{promoted\_week}=\varnothing$). Each condition encodes a property a genuine theme
must have over time.

\begin{description}[leftmargin=1.6em,style=nextline]
  \item[Persistence \eqref{eq:persist}.] Requiring survival for at least $\rho=2$ caught weeks
        rejects one-week flashes. A spike that appears and vanishes is an event; a theme is a
        process with positive duration. This is a minimum-lifetime (survival) filter on the
        anchor's caught history.
  \item[Growth \eqref{eq:grow}.] $d_i\ge\max_{j<i}d_j$ asks that week $w_i$ be a \emph{record}
        of the reach sequence: the theme is recruiting strictly more (or at least as many)
        distinct partners than in any prior week. A maturing theme expands its cross-sectional
        footprint, so its reach sequence should keep setting new highs; a saturated or fading
        story does not. Demanding a running-maximum, rather than a mere positive slope, makes
        the test robust to week-to-week noise.
  \item[Bridge \eqref{eq:bridge}.] The cumulative \emph{median} clustering must stay at or below
        $\kappa=0.60$, so the \emph{typical} week is a low-clustering hub. The median (rather
        than the mean or the latest value) is used because a real theme can contain a transient
        sub-event that briefly spikes clustering; the median ignores such outliers and asks only
        that the anchor bridges unrelated stories in most of its life.
\end{description}

\begin{remark}[Catch vs.\ promote = recall vs.\ precision]
Catch thresholds the static features $(c_w,\deg_w)$ of a single week; promote thresholds the
\emph{dynamics} $(d_i)_i$ and $(\clu_i)_i$ across weeks. Splitting the work this way keeps the
candidate set large until the cheap temporal evidence (persistence, growth, low median
clustering) has accumulated, then prunes hard. On the benchmark, catch keeps $814$ entities and
promote keeps $20$.
\end{remark}

% =====================================================================
\section{Stage 4: assembling themes}\label{sec:stage4}
% =====================================================================

A single theme often surfaces under several anchor variants (e.g.\ \texttt{chatgpt} and
\texttt{chatgpt maker}). Let $S$ be the set of promoted anchors and build the graph
\begin{equation}
\mathcal{H}=(S,\,F),\qquad
(a,b)\in F \iff a,b\in S,\ a\ne b,\ b\in\mathrm{subgraph}(a),
\end{equation}
where $\mathrm{subgraph}(a)$ is the anchor plus its top entity partners recorded in $\Rtab$.
The \emph{themes} are the connected components of $\mathcal{H}$:
\begin{equation}
\text{themes}=\bigl\{\,\text{connected components of }\mathcal{H}\,\bigr\}.
\end{equation}
Mathematically, the themes are the equivalence classes of the transitive closure of the
co-occurrence relation on promoted anchors. This deduplicates variant names into one object and
merges promoted entities that belong to the same narrative, taking $20$ promoted anchors to
$18$ themes. Each theme is summarised by its merged subgraph and its reach
$\max_{a}\deg_{\max}(a)$.

% =====================================================================
\section{Stage 5: the language-model reject filter}\label{sec:stage5}
% =====================================================================

The surviving themes are passed, top-by-reach, to a language model acting as a
\emph{conservative, reject-only} classifier. It never decides what \emph{is} a theme; it only
flags a cluster for removal when, judging from the headlines alone, it clearly falls into one of
\begin{enumerate}[leftmargin=1.6em,nosep]
  \item \texttt{single\_entity\_event} -- one firm's idiosyncratic event, no multi-actor narrative;
  \item \texttt{macro\_market\_aggregate} -- rates, inflation, FX, indices, central-bank policy, commodities;
  \item \texttt{boilerplate\_wire} -- wire formatting, calendars, routine ``shares rise/fall'' PR.
\end{enumerate}
The rule is \emph{when in doubt, keep}. Writing $K(g)\in\{0,1\}$ for the keep decision on a
group $g$, the final shortlist is $\{g:\,K(g)=1\}$ ($4$ of $18$ on the benchmark, including the
generative-AI theme).

\paragraph{Why a reject-only, keep-biased filter at the end.}
The pipeline is a \emph{cascade}: cheap, recall-oriented graph tests at the front, an expensive
but accurate test at the back, applied only to the few survivors. Two properties make the
ordering correct.
\begin{itemize}[leftmargin=1.6em]
  \item \textbf{Monotone precision.} Because the model can only \emph{remove} groups, it cannot
        introduce false positives. It strictly raises precision and leaves recall bounded below
        by the recall of Stages 0--4. The ``keep when unsure'' bias means its own false-rejection
        rate is small, so little true recall is lost.
  \item \textbf{Cost placement.} Let the front stages produce $n$ candidates at negligible cost
        and the LLM cost $C_{\text{LLM}}$ per group with error rate $\varepsilon$. Applying the
        LLM only to the $|{\text{themes}}|\ll n$ survivors makes total cost
        $O(|{\text{themes}}|\,C_{\text{LLM}})$ while the residual false-positive rate is the LLM's
        own $\varepsilon$. Running the expensive judge first, or as a detector rather than a
        filter, would multiply cost by orders of magnitude for no recall gain.
\end{itemize}
In Bayesian terms, the graph features supply a high-recall prior over candidate themes and the
LLM supplies a likelihood-ratio test on the assembled headline evidence, evaluated only where the
prior is already high.

% =====================================================================
\section{Why the whole pipeline works}
% =====================================================================

Two robust, unsupervised signals carry the result and a cheap selector plus an LLM clean its
ends.
\begin{enumerate}[leftmargin=1.6em]
  \item \textbf{Novelty} (Stage 1) guarantees the detector looks at \emph{emerging} vocabulary
        only, by set difference against history.
  \item \textbf{Graph topology} (Stage 2--3) distinguishes a \emph{theme} (high-reach,
        low-clustering hub spanning structural holes) from an \emph{event} (high-clustering
        clique) and from \emph{noise} (low reach), using only co-occurrence counts.
  \item \textbf{Temporal dynamics} (Stage 3) demand the structural signature \emph{persist},
        \emph{grow}, and \emph{stay a bridge}, which is what separates a durable theme from a
        flash.
  \item \textbf{Entity selection} (Stage 0) and the \textbf{LLM filter} (Stage 5) clean the two
        ends of the funnel; neither is what makes detection work, which is why swapping the
        hand-written blocklist for an NER vocabulary changes the entry/exit weeks slightly but
        reproduces the same outcome.
\end{enumerate}

\paragraph{Benchmark.}
On the generative-AI theme the funnel runs
$3{,}767\to 814\to 20\to 18\to 4$, with \texttt{chatgpt} first caught in the week of
3--9 January 2023 and promoted in the week of 17--23 January 2023. Against a thematic-ETF
inception of 17 May 2023 this is a lead of about $120$ days ($\approx 4$ months), and the LLM
keeps the cluster. The result reproduces the predecessor's manual-blocklist run, confirming that
the detection is driven by the novelty-and-topology core rather than by either Stage-0 device.

% =====================================================================
\section*{Appendix: parameters}
% =====================================================================
\addcontentsline{toc}{section}{Appendix: parameters}

\begin{center}
\begin{tabular}{lll}
\toprule
Symbol & Value & Meaning \\
\midrule
$t_0$              & 2022-09-30 & baseline cutoff \\
$[t_1,t_2]$        & 2022-10-01 .. 2023-02-28 & discovery window \\
freq               & \texttt{W-MON} & weekly bucketing \\
$m$                & $3$  & minimum weekly mentions (support) \\
$\deg_{\min}$      & $8$  & catch reach threshold \\
$K$                & $15$ & partners used for clustering \\
$\rho$             & $2$  & persistence (weeks) \\
$\kappa$           & $0.60$ & maximum median clustering (bridge) \\
$\mathcal{L}$      & ORG, PRODUCT, PERSON, \dots & NER labels kept for $\Ent$ \\
NER vocab          & $18{,}314$ tokens & non-dictionary entity tokens, count $\ge 2$ \\
\bottomrule
\end{tabular}
\end{center}

\end{document}
